\chapter{Three-Way Comparison --- The Spaceship Operator}
\label{ch:c20-spaceship}

\begin{quote}
\emph{C++20's three-way comparison operator \texttt{<=>} --- the ``spaceship'' --- replaces the six hand-written comparison operators most types needed with a single defaulted declaration, and gives the language a principled vocabulary for the three kinds of ordering a type can have. This chapter covers the operator, its ordering category types, defaulted versus custom comparisons, the operator-rewriting rules, and the performance subtleties.}
\end{quote}

The spaceship operator collapses six operators to one function whose return value encodes the full ordering; the compiler rewrites the relational operators in terms of it.

\section{The Problem: Six Operators Per Type}
\label{sec:c20-six-operators}

\begin{lstlisting}[language=C++,caption={The pre-C++20 boilerplate --- six operators by hand}]
struct Version {
    int major, minor, patch;
    bool operator==(const Version& o) const {
        return major==o.major && minor==o.minor && patch==o.patch;
    }
    bool operator!=(const Version& o) const { return !(*this == o); }
    bool operator<(const Version& o) const {
        if (major != o.major) return major < o.major;
        if (minor != o.minor) return minor < o.minor;
        return patch < o.patch;
    }
    bool operator>(const Version& o) const  { return o < *this; }
    bool operator<=(const Version& o) const { return !(o < *this); }
    bool operator>=(const Version& o) const { return !(*this < o); }
};
\end{lstlisting}

\section{The Spaceship Operator and Defaulted Comparison}
\label{sec:c20-defaulted-comparison}

\begin{lstlisting}[language=C++,caption={The entire block replaced by two defaulted lines}]
#include <compare>

struct Version {
    int major, minor, patch;
    auto operator<=>(const Version&) const = default;   // generates <, >, <=, >=
    bool operator==(const Version&) const = default;    // generates ==, !=
};
\end{lstlisting}

\texttt{= default} on \texttt{<=>} produces a member-wise, in-declaration-order, lexicographic comparison. \texttt{<compare>} is required.

\section{The Three Ordering Categories}
\label{sec:c20-ordering-categories}

\begin{center}
\begin{tabular}{lll}
\hline
\textbf{Return type} & \textbf{Meaning} & \textbf{Key property} \\
\hline
\texttt{strong\_ordering} & total order, substitutable & \texttt{a==b} $\Rightarrow$ indistinguishable \\
\texttt{weak\_ordering} & total order, not substitutable & equivalence classes \\
\texttt{partial\_ordering} & some pairs unordered & \texttt{NaN <=> x} is \texttt{unordered} \\
\hline
\end{tabular}
\end{center}

\begin{lstlisting}[language=C++,caption={The named result values of each category}]
#include <compare>

void inspect(std::strong_ordering o) {
    if (o == std::strong_ordering::less)    { /* a < b */ }
    if (o == std::strong_ordering::equal)   { /* fully substitutable */ }
    if (o == std::strong_ordering::greater) { /* a > b */ }
}
void inspect_fp(std::partial_ordering o) {
    if (o == std::partial_ordering::unordered) { /* involves NaN */ }
}
\end{lstlisting}

\texttt{strong} has \texttt{less/equal/greater}; \texttt{weak} has \texttt{less/equivalent/greater}; \texttt{partial} adds \texttt{unordered}. Integers yield strong; \texttt{float}/\texttt{double} yield partial.

\section{Operator Rewriting and Synthesized Candidates}
\label{sec:c20-operator-rewriting}

\begin{lstlisting}[language=C++,caption={How the compiler rewrites relational operators}]
// Source you write:          Compiler rewrites to:
//   a < b           ->          (a <=> b) < 0
//   a > b           ->          (a <=> b) > 0
//   a <= b          ->          (a <=> b) <= 0
//   a >= b          ->          (a <=> b) >= 0
//   a > b           ->          0 < (b <=> a)    // synthesized (reversed) candidate
\end{lstlisting}

A single \texttt{friend auto operator<=>} makes all eight relational comparisons (both orders) work. The rewrite applies only to \texttt{<,>,<=,>=}.

\section{== Is Special: Why Equality Is Separate}
\label{sec:c20-equality-special}

\begin{lstlisting}[language=C++,caption={Equality must be defaulted separately}]
#include <compare>
#include <string>
#include <vector>

struct Record {
    int id;
    std::vector<std::string> tags;
    auto operator<=>(const Record&) const = default;   // lexicographic ordering
    bool operator==(const Record&) const = default;    // cheaper: short-circuits on size
};
\end{lstlisting}

\texttt{==} can short-circuit on size in O(1); \texttt{<=>} must compare element-by-element. \texttt{!=} is always rewritten from \texttt{==}.

\section{Custom Three-Way Comparisons}
\label{sec:c20-custom-spaceship}

\begin{lstlisting}[language=C++,caption={A custom <=> with a chosen ordering category}]
#include <compare>
#include <string>
#include <cctype>

struct CaseInsensitive {
    std::string s;
    std::weak_ordering operator<=>(const CaseInsensitive& o) const {
        auto lower = [](std::string x){
            for (char& c : x) c = static_cast<char>(std::tolower(c));
            return x;
        };
        return lower(s) <=> lower(o.s);   // narrow string's strong order to weak
    }
    bool operator==(const CaseInsensitive& o) const { return (*this <=> o) == 0; }
};
\end{lstlisting}

Standard helpers: \texttt{std::compare\_three\_way}, and \texttt{std::strong\_order}/\texttt{weak\_order}/\texttt{partial\_order}.

\section{Member-wise Semantics and Ordering Strength Deduction}
\label{sec:c20-strength-deduction}

\begin{lstlisting}[language=C++,caption={Deduced category is the weakest among members}]
#include <compare>

struct Mixed {
    int    n;     // contributes strong_ordering
    double d;     // contributes partial_ordering
    auto operator<=>(const Mixed&) const = default;
    // deduced return type: std::partial_ordering (weakest member wins)
};
\end{lstlisting}

\texttt{std::common\_comparison\_category} computes the result. Name the category explicitly to assert a minimum strength.

\section{Performance Considerations}
\label{sec:c20-spaceship-perf}

\begin{itemize}
\item \textbf{Prefer \texttt{==} for equality on containers} --- it short-circuits on size; \texttt{<=>} does not.
\item \textbf{A single \texttt{<=>} call yields all relations} --- capture \texttt{auto c = a <=> b;} once in hot comparators.
\item \textbf{\texttt{partial\_ordering} carries an extra state} --- four outcomes, marginally heavier; NaN keys make ordered containers ill-formed.
\item \textbf{Defaulted comparisons are \texttt{constexpr}-friendly} when members are.
\end{itemize}

\section{Professional Insights}
\label{sec:c20-spaceship-insights}

\textbf{Default both \texttt{<=>} and \texttt{==}.} The review-friendly habit; guarantees the fast \texttt{==} path for container members.

\textbf{Choose the ordering category deliberately --- it is part of the contract.} A case-insensitive string returning \texttt{strong\_ordering} is a latent bug; return \texttt{weak\_ordering}.

\textbf{Compute \texttt{<=>} once and reuse it in comparison-heavy code} --- halves work in sort comparators and tree inserts.

\textbf{Floating-point keys yield \texttt{partial\_ordering}, and NaN breaks ordered containers} --- use \texttt{std::strong\_order} or guarantee no NaNs.

\textbf{Reach for \texttt{<=>} to delete CRTP comparison helpers} --- removes boilerplate, shrinks inline footprint, eliminates inconsistency bugs.
